\documentclass[12pt,a4paper]{article}

% ─── Paquetes ─────────────────────────────────────────────────────────────────
\usepackage[utf8]{inputenc}
\usepackage[T1]{fontenc}
\usepackage[spanish]{babel}
\usepackage{geometry}
\geometry{margin=2.5cm}

\usepackage{tabularx}
\usepackage{booktabs}
\usepackage{longtable}
\usepackage{array}
\usepackage{xcolor}
\usepackage{hyperref}
\usepackage{listings}
\usepackage{caption}
\usepackage{enumitem}
\usepackage{graphicx}
\usepackage{fancyhdr}
\usepackage{titlesec}
\usepackage{microtype}
\usepackage{parskip}

% ─── Configuración de Hyperref ──────────────────────────────────────────────
\hypersetup{
    colorlinks=true,
    linkcolor=blue!60!black,
    urlcolor=blue!60!black,
    citecolor=green!40!black,
    pdftitle={Especificación Técnica: Lógica de Perfilado y Detección Reactiva en EdgeLeak AI},
    pdfauthor={EdgeLeak AI — Proyecto PICUR Uniremington},
    pdfsubject={Documentación de Arquitectura IoT},
    pdfkeywords={IoT, Flutter, ESP32, detección de fugas, machine learning, SQLite},
}

% ─── Configuración de Listings ──────────────────────────────────────────────
\definecolor{codegreen}{rgb}{0,0.5,0}
\definecolor{codegray}{rgb}{0.5,0.5,0.5}
\definecolor{codeblue}{rgb}{0.1,0.2,0.7}
\definecolor{backcolour}{rgb}{0.97,0.97,0.97}

\lstdefinestyle{dartstyle}{
    backgroundcolor=\color{backcolour},
    commentstyle=\color{codegreen}\itshape,
    keywordstyle=\color{codeblue}\bfseries,
    numberstyle=\tiny\color{codegray},
    stringstyle=\color{codegreen!60!black},
    basicstyle=\ttfamily\footnotesize,
    breakatwhitespace=false,
    breaklines=true,
    captionpos=b,
    keepspaces=true,
    numbers=left,
    numbersep=5pt,
    showspaces=false,
    showstringspaces=false,
    showtabs=false,
    tabsize=2,
    frame=single,
    rulecolor=\color{gray!30},
}
\lstset{style=dartstyle}

% ─── Encabezados y Pies de Página ──────────────────────────────────────────
\pagestyle{fancy}
\fancyhf{}
\fancyhead[L]{\small EdgeLeak AI — Documentación Técnica v2.0}
\fancyhead[R]{\small Proyecto PICUR Uniremington}
\fancyfoot[C]{\thepage}
\renewcommand{\headrulewidth}{0.4pt}

% ─── Configuración de Secciones ─────────────────────────────────────────────
\titleformat{\section}[block]{\large\bfseries\color{blue!70!black}}{\thesection.}{0.5em}{}[\titlerule]
\titleformat{\subsection}[block]{\normalsize\bfseries\color{blue!50!black}}{\thesubsection.}{0.5em}{}
\titleformat{\subsubsection}[block]{\normalsize\itshape}{\thesubsubsection.}{0.5em}{}

% ─── Colores del sistema ─────────────────────────────────────────────────────
\definecolor{estadoNormal}{RGB}{76,175,80}
\definecolor{estadoAnomalia}{RGB}{255,193,7}
\definecolor{estadoFuga}{RGB}{244,67,54}

% ─── Documento ───────────────────────────────────────────────────────────────

\begin{document}

% ── Portada ──────────────────────────────────────────────────────────────────
\begin{titlepage}
    \centering
    \vspace*{2cm}
    {\Huge\bfseries EdgeLeak AI\par}
    \vspace{0.5cm}
    {\Large\itshape Especificación Técnica\par}
    \vspace{1cm}
    \rule{\linewidth}{0.5pt}
    \vspace{0.5cm}
    {\large\bfseries Lógica de Perfilado de Usuario y Detección\\[4pt]
    Reactiva de Microfugas\par}
    \vspace{0.5cm}
    \rule{\linewidth}{0.5pt}
    \vspace{2cm}

    \begin{tabularx}{0.85\textwidth}{lX}
        \toprule
        \textbf{Proyecto} & PICUR — Uniremington \\
        \textbf{Versión del documento} & 2.0 \\
        \textbf{Clasificación} & Documento técnico de arquitectura \\
        \textbf{Sistema objetivo} & EdgeLeak AI — Detección de microfugas en lavaplatos \\
        \textbf{Hardware} & ESP32 + YF-S201 + KY-037 \\
        \textbf{Plataforma de software} & Flutter 3.x / Dart 3.x / SQLite \\
        \textbf{Servicio de IA} & Groq API (Llama 3.3 70B) \\
        \bottomrule
    \end{tabularx}

    \vfill
    {\small Sistema de Detección de Microfugas para Lavaplatos Manuales\\
    Arquitectura Edge-to-Cloud con Inteligencia Artificial Contextual\par}
\end{titlepage}

% ── Tabla de Contenido ───────────────────────────────────────────────────────
\tableofcontents
\newpage

% ─────────────────────────────────────────────────────────────────────────────
\section{Introducción y Alcance}
% ─────────────────────────────────────────────────────────────────────────────

El sistema \textbf{EdgeLeak AI} es una solución de Internet de las Cosas (IoT) orientada a la detección de microfugas en lavaplatos manuales domésticos e industriales. Su arquitectura combina procesamiento en el borde (\textit{edge computing}) con análisis contextual en la nube mediante la API de Groq, siguiendo el paradigma \textit{Edge-to-Cloud}.

El presente documento especifica formalmente la lógica de perfilado de usuario implementada en la versión 2.0, incluyendo:

\begin{itemize}[noitemsep]
    \item La máquina de estados de tres niveles (\texttt{Normal}, \texttt{Anomalía}, \texttt{Fuga}).
    \item El filtro digital de ruido implementado en el servicio \texttt{RuidoFilterService}.
    \item El módulo de perfilado adaptativo \texttt{BaselineService} con agregación de datos jerárquica.
    \item El modelo de invocación reactiva (event-driven) de la IA de Groq.
    \item El esquema de datos de la base de datos local SQLite versión 5.
\end{itemize}

\subsection{Objetivos del Sistema}

\begin{enumerate}[noitemsep]
    \item Detectar con exactitud tres niveles operacionales del sistema hídrico monitoreado.
    \item Distinguir entre impactos acústicos aislados y patrones rítmicos de goteo.
    \item Aprender el comportamiento normal del usuario en bloques de 3 a 5 días (\textit{baseline}).
    \item Invocar la IA de Groq exclusivamente ante eventos de desviación significativa del baseline.
    \item Agregar los datos crudos en resúmenes jerárquicos para optimizar el almacenamiento SQLite.
    \item Notificar al usuario mediante correo electrónico ante cualquier cambio de estado no-normal.
\end{enumerate}

% ─────────────────────────────────────────────────────────────────────────────
\section{Arquitectura Edge-to-Cloud}
% ─────────────────────────────────────────────────────────────────────────────

El diseño arquitectónico de EdgeLeak AI sigue tres capas funcionales:

\subsection{Capa Edge (ESP32 + Sensores)}

El microcontrolador ESP32 recopila datos de dos sensores cada 5 segundos y los transmite al servidor HTTP local de la aplicación Flutter:

\begin{itemize}[noitemsep]
    \item \textbf{YF-S201}: Sensor de flujo de agua. Genera pulsos proporcionales al caudal. La unidad de medida procesada es litros por minuto (L/min).
    \item \textbf{KY-037}: Micrófono de alta sensibilidad. Entrega un valor analógico-digital (ADC) proporcional al nivel de presión acústica ambiente.
\end{itemize}

La comunicación se realiza mediante peticiones HTTP POST al endpoint \texttt{/api/sensor} del servidor local (puerto 8080), protegido mediante cabecera \texttt{x-api-key}.

\subsection{Capa de Procesamiento Local (Flutter / Dart)}

La aplicación Flutter ejecuta el procesamiento principal en el hilo de la interfaz de usuario (UI thread) con operaciones asíncronas que no bloquean la renderización:

\begin{enumerate}[noitemsep]
    \item \textbf{Recepción}: \texttt{LocalApiService} recibe y valida la petición del ESP32.
    \item \textbf{Clasificación}: \texttt{DashboardController} aplica la máquina de estados de 3 niveles.
    \item \textbf{Filtrado de ruido}: \texttt{RuidoFilterService} distingue impactos aislados de patrones de goteo.
    \item \textbf{Detección de desviación}: \texttt{BaselineService} compara contra el perfil aprendido.
    \item \textbf{Persistencia}: \texttt{DatabaseService} almacena lecturas crudas y agrega en resúmenes.
    \item \textbf{Notificación}: \texttt{EmailService} envía alerta al correo del usuario registrado.
\end{enumerate}

\subsection{Capa Cloud (Groq API — Llama 3.3 70B)}

La IA de Groq actúa como árbitro final de clasificación y es invocada de forma reactiva. El payload incluye la lectura anómala actual más el contexto histórico del usuario (resúmenes de 5 días y el historial mensual más reciente). La respuesta determina el veredicto definitivo entre los tres estados canónicos del sistema.

% ─────────────────────────────────────────────────────────────────────────────
\section{Máquina de Estados de Tres Niveles}
% ─────────────────────────────────────────────────────────────────────────────

\subsection{Definición Formal}

El estado del sistema se representa mediante la enumeración \texttt{EstadoSensor}, que define tres valores canónicos mutuamente excluyentes:

\begin{center}
\begin{tabularx}{0.9\textwidth}{>{\bfseries}lXX}
    \toprule
    \textbf{Estado} & \textbf{Condición de activación} & \textbf{Acción del sistema} \\
    \midrule
    \textcolor{estadoNormal}{Normal} &
        Flujo $\leq 0{,}5$ L/min \textit{y} sin patrón de goteo detectado por el filtro de ruido. &
        Monitoreo pasivo. Sin notificaciones. \\
    \addlinespace
    \textcolor{estadoAnomalia}{Anomalía} &
        Flujo $\in (0{,}5; 5{,}0]$ L/min, \textit{o} patrón de goteo rítmico detectado por el filtro, \textit{o} lógica nocturna activa (flujo = 0 + ruido constante $>$ base durante 0:00--6:00). &
        Alerta inmediata por correo. Invocación reactiva de la IA tras 3 strikes. \\
    \addlinespace
    \textcolor{estadoFuga}{Fuga} &
        Flujo $> 5{,}0$ L/min (umbral crítico de caudal activo). &
        Alerta crítica inmediata por correo. Invocación directa de la IA sin esperar strikes. \\
    \bottomrule
\end{tabularx}
\captionof{table}{Tabla de condiciones de transición de la máquina de estados.}
\end{center}

\subsection{Jerarquía de Prioridad}

La clasificación respeta una jerarquía estricta de evaluación secuencial:

\begin{enumerate}[noitemsep]
    \item \textbf{Fuga} (prioridad máxima): Evaluada primero. Si el caudal supera 5,0 L/min, el estado se establece como \texttt{Fuga} sin importar el nivel de ruido.
    \item \textbf{Anomalía}: Evaluada si no se cumplió la condición de Fuga. Requiere flujo intermedio \textit{o} patrón de goteo confirmado por \texttt{RuidoFilterService}.
    \item \textbf{Normal} (prioridad mínima): Estado por defecto cuando ninguna de las condiciones anteriores se cumple.
\end{enumerate}

\subsection{Reactividad de la Interfaz de Usuario}

La UI reacciona de forma síncrona al cambio de estado gracias al patrón \texttt{ChangeNotifier} de Flutter. La actualización visual es inmediata (antes de las operaciones de base de datos o la llamada a la IA):

\begin{center}
\begin{tabularx}{0.88\textwidth}{>{\bfseries}lXXX}
    \toprule
    \textbf{Estado} & \textbf{Color dominante} & \textbf{Modo de olas} & \textbf{Velocidad de animación} \\
    \midrule
    Normal   & Azul (\texttt{\#2196F3}) & Normal  & 2000 ms/ciclo \\
    Anomalía & Amarillo (\texttt{\#FFC107}) & Anomalia & 1000 ms/ciclo \\
    Fuga     & Rojo (\texttt{\#F44336})  & Fuga    & 700 ms/ciclo \\
    \bottomrule
\end{tabularx}
\captionof{table}{Mapeado de estados a parámetros visuales de la UI.}
\end{center}

% ─────────────────────────────────────────────────────────────────────────────
\section{Filtro Digital del Sensor de Ruido (KY-037)}
% ─────────────────────────────────────────────────────────────────────────────

\subsection{Problema Técnico}

El sensor KY-037 es susceptible a falsas alarmas generadas por impactos mecánicos aislados (por ejemplo, la caída de un utensilio metálico), los cuales producen picos de amplitud ADC similares o superiores a los patrones de goteo real. Sin un mecanismo de filtrado, cada impacto aislado clasificaría erróneamente el sistema como \texttt{Anomalía}.

\subsection{Algoritmo de Ventana Deslizante}

El servicio \texttt{RuidoFilterService} implementa un filtro de ventana deslizante temporal con las siguientes características:

\begin{itemize}[noitemsep]
    \item \textbf{Ventana máxima}: 12 muestras o 60 segundos (lo que se alcance primero).
    \item \textbf{Criterio de patrón rítmico}: Al menos 3 muestras dentro de la ventana con valor ADC superior al umbral de ruido significativo (\texttt{\_umbralRuido = 1500 ADC}).
    \item \textbf{Criterio de impacto aislado}: Exactamente 1 ó 2 muestras sobre el umbral, rodeadas de muestras en silencio.
\end{itemize}

El resultado se devuelve como un valor de la enumeración \texttt{ResultadoRuido}:

\begin{center}
\begin{tabularx}{0.88\textwidth}{>{\ttfamily}lX}
    \toprule
    \textbf{Valor} & \textbf{Semántica} \\
    \midrule
    silencio       & Ninguna muestra sobre el umbral. Sin acción. \\
    impactoAislado & Pico único. Se ignora como posible fuga. \\
    patronGoteo    & Patrón rítmico sostenido. Clasifica como \texttt{Anomalía}. \\
    \bottomrule
\end{tabularx}
\captionof{table}{Valores de salida de \texttt{RuidoFilterService.clasificar()}.}
\end{center}

\subsection{Lógica Nocturna}

La lógica nocturna cubre el escenario de goteo en la boquilla del lavaplatos durante horas de inactividad (0:00--6:00 a.m.), donde el sensor de flujo no registra caudal porque el goteo no es suficiente para activar el contador del YF-S201.

\textbf{Condiciones de activación simultáneas}:
\begin{enumerate}[noitemsep]
    \item Hora del sistema comprendida en el intervalo $[0, 6)$.
    \item Flujo reportado por el ESP32 igual a 0,0 L/min.
    \item Promedio del nivel de ruido en la ventana actual superior al umbral de silencio base (\texttt{\_umbralSilencioBase = 800 ADC}).
    \item Mínimo 3 muestras presentes en la ventana (validez estadística).
\end{enumerate}

Si todas las condiciones se satisfacen simultáneamente, el servicio devuelve \texttt{ResultadoRuido.patronGoteo}, lo que eleva el estado del sistema a \texttt{Anomalía}.

\subsection{Tabla de Umbrales de Calibración}

\begin{center}
\begin{tabularx}{0.95\textwidth}{lXX>{\ttfamily}l}
    \toprule
    \textbf{Parámetro} & \textbf{Descripción} & \textbf{Criterio técnico} & \textbf{Valor por defecto} \\
    \midrule
    Umbral de ruido significativo &
        Nivel ADC mínimo para considerar actividad acústica relevante. &
        Calibrado sobre el nivel de ruido de flujo normal activo. &
        1500 \\
    \addlinespace
    Umbral de silencio base &
        Nivel ADC máximo durante períodos sin actividad acústica. &
        Medición en ambiente sin agua durante la noche. &
        800 \\
    \addlinespace
    Mínimo de muestras de goteo &
        Cantidad mínima de muestras sobre el umbral para confirmar patrón rítmico. &
        Evita falsas alarmas por impactos dobles casuales. &
        3 \\
    \addlinespace
    Tamaño de ventana (muestras) &
        Capacidad máxima del buffer circular. &
        Cubre aproximadamente 1 minuto a 5 s/lectura. &
        12 \\
    \addlinespace
    Duración de ventana temporal &
        Tiempo máximo de antigüedad de una muestra para permanecer en la ventana. &
        Las muestras antiguas no contribuyen al análisis de patrón actual. &
        60 s \\
    \addlinespace
    Hora de inicio nocturno &
        Inicio del período de lógica nocturna (hora del dispositivo). &
        Horas de mínima actividad del usuario. &
        0:00 \\
    \addlinespace
    Hora de fin nocturno &
        Fin del período de lógica nocturna (exclusive). &
        A partir de las 6:00, el usuario puede comenzar a usar el lavaplatos. &
        6:00 \\
    \bottomrule
\end{tabularx}
\captionof{table}{Tabla de Umbrales de Calibración del Filtro de Ruido KY-037.}
\end{center}

% ─────────────────────────────────────────────────────────────────────────────
\section{Perfilado de Usuario y Gestión del Baseline}
% ─────────────────────────────────────────────────────────────────────────────

\subsection{Descripción del Baseline}

El \textit{baseline} (línea base) representa el perfil estadístico del comportamiento hídrico y acústico normal del usuario, construido a partir de las lecturas crudas almacenadas en la tabla \texttt{lecturas\_raw} de los últimos \texttt{\_diasBaseline = 5} días.

El baseline se segmenta en dos bloques horarios para capturar la diferencia natural entre el uso diurno y el ambiente nocturno:

\begin{itemize}[noitemsep]
    \item \textbf{Período diurno}: 6:00--22:00. El usuario realiza usos activos del lavaplatos.
    \item \textbf{Período nocturno}: 22:00--6:00. Silencio esperado; solo ruido ambiental de fondo.
\end{itemize}

El \texttt{BaselineService} calcula el promedio de flujo y ruido de cada período por separado. Al recibir una nueva lectura, selecciona el período correspondiente a la hora actual para comparar.

\subsection{Proceso de Construcción}

El baseline se reconstruye automáticamente:
\begin{itemize}[noitemsep]
    \item Al iniciar la aplicación.
    \item Al cambiar de usuario (llamada explícita a \texttt{reconstruirBaseline()}).
    \item Tras cada ciclo de agregación de 24 horas.
\end{itemize}

Se requiere un mínimo de 30 lecturas crudas para que el baseline sea estadísticamente válido. Con datos insuficientes, el módulo de detección de desviación retorna \texttt{true} de forma conservadora (se escala el procesamiento al modo reactivo completo).

\subsection{Detección de Desviación Significativa}

La función \texttt{esDesviacionSignificativa(ruido, flujo)} compara la lectura actual contra los promedios del baseline del período activo:

\begin{equation}
    \text{desvíaFlujo} = |flujo_{\text{actual}} - \overline{flujo}_{\text{baseline}}| > \Delta F_{\text{umbral}}
\end{equation}
\begin{equation}
    \text{desvíaRuido} = |ruido_{\text{actual}} - \overline{ruido}_{\text{baseline}}| > \Delta R_{\text{umbral}}
\end{equation}
\begin{equation}
    \text{esDesviación} = \text{desvíaFlujo} \lor \text{desvíaRuido}
\end{equation}

Donde $\Delta F_{\text{umbral}} = 0{,}30$ L/min y $\Delta R_{\text{umbral}} = 500$ ADC son los umbrales de desviación mínimos.

Solo cuando \texttt{esDesviación = true} \textit{y} el contador de strikes alcanza el umbral, se invoca la API de Groq, eliminando el polling innecesario.

% ─────────────────────────────────────────────────────────────────────────────
\section{Modelo de Datos y Agregación Jerárquica (SQLite v5)}
% ─────────────────────────────────────────────────────────────────────────────

\subsection{Esquema de Base de Datos}

La base de datos local \texttt{edgeleak\_v5.db} (versión 5) contiene cinco tablas:

\begin{center}
\begin{tabularx}{0.95\textwidth}{lX}
    \toprule
    \textbf{Tabla} & \textbf{Propósito} \\
    \midrule
    \texttt{usuarios}            & Registro de cuentas de usuario con roles (operador/administrador). \\
    \texttt{historial}           & Alertas confirmadas por la IA, con veredicto, severidad y mensaje. \\
    \texttt{lecturas\_raw}       & Lecturas crudas individuales del ESP32 (ruido ADC + flujo L/min). \\
    \texttt{lecturas\_resumen\_5d} & Resúmenes agregados de bloques de 5 días de lecturas crudas. \\
    \texttt{historial\_mensual}  & Consolidación de múltiples resúmenes de 5 días en un registro mensual. \\
    \bottomrule
\end{tabularx}
\captionof{table}{Tablas de la base de datos \texttt{edgeleak\_v5.db}.}
\end{center}

\subsection{Modelo de Datos de Agregación}

\begin{center}
\begin{tabularx}{0.95\textwidth}{llXl}
    \toprule
    \multicolumn{4}{c}{\textbf{Tabla: \texttt{lecturas\_resumen\_5d}}} \\
    \midrule
    \textbf{Columna} & \textbf{Tipo} & \textbf{Descripción} & \textbf{Restricción} \\
    \midrule
    \texttt{id}                  & INTEGER & Identificador autoincremental.                   & PK \\
    \texttt{fecha\_inicio}       & TEXT    & Inicio del bloque de 5 días (ISO 8601).          & NOT NULL \\
    \texttt{fecha\_fin}          & TEXT    & Fin del bloque de 5 días (ISO 8601).             & NOT NULL \\
    \texttt{flujo\_promedio}     & REAL    & Promedio aritmético del caudal en el período.    & NOT NULL \\
    \texttt{ruido\_promedio}     & REAL    & Promedio aritmético del ruido ADC en el período. & NOT NULL \\
    \texttt{total\_lecturas}     & INTEGER & Cantidad de lecturas crudas agregadas.            & NOT NULL \\
    \texttt{estado\_predominante}& TEXT    & Estado con mayor frecuencia en el período.        & NOT NULL \\
    \bottomrule
\end{tabularx}
\vspace{0.5em}

\begin{tabularx}{0.95\textwidth}{llXl}
    \toprule
    \multicolumn{4}{c}{\textbf{Tabla: \texttt{historial\_mensual}}} \\
    \midrule
    \textbf{Columna} & \textbf{Tipo} & \textbf{Descripción} & \textbf{Restricción} \\
    \midrule
    \texttt{id}              & INTEGER & Identificador autoincremental.                        & PK \\
    \texttt{fecha\_inicio}   & TEXT    & Inicio del período mensual consolidado.               & NOT NULL \\
    \texttt{fecha\_fin}      & TEXT    & Fin del período mensual consolidado.                  & NOT NULL \\
    \texttt{num\_anomalias}  & INTEGER & Número de bloques de 5 días clasificados como Anomalía. & NOT NULL \\
    \texttt{num\_fugas}      & INTEGER & Número de bloques de 5 días clasificados como Fuga.    & NOT NULL \\
    \texttt{flujo\_promedio} & REAL    & Promedio de flujo a lo largo del período mensual.     & NOT NULL \\
    \texttt{ruido\_promedio} & REAL    & Promedio de ruido a lo largo del período mensual.     & NOT NULL \\
    \texttt{resumen\_json}   & TEXT    & JSON con los bloques de 5 días que componen el mes.   & NOT NULL \\
    \bottomrule
\end{tabularx}
\captionof{table}{Modelos de datos de la agregación jerárquica.}
\end{center}

\subsection{Proceso de Agregación en Segundo Plano}

El \texttt{BaselineService} ejecuta un timer periódico cada 24 horas con la siguiente cadena de operaciones:

\begin{enumerate}[noitemsep]
    \item \textbf{Agregación de lecturas crudas}: Se calculan el promedio de flujo, el promedio de ruido, el total de lecturas y el estado predominante del bloque de 5 días. El resultado se inserta en \texttt{lecturas\_resumen\_5d}.
    \item \textbf{Reconstrucción del baseline}: Se recalcula el perfil estadístico del usuario con los datos más recientes.
    \item \textbf{Consolidación mensual}: Si existen 5 o más resúmenes de 5 días, se genera un registro en \texttt{historial\_mensual} y los resúmenes individuales correspondientes se eliminan de la tabla \texttt{lecturas\_resumen\_5d} para liberar espacio.
\end{enumerate}

Este proceso es transparente para el usuario: no genera ninguna notificación ni bloquea la interfaz de usuario.

% ─────────────────────────────────────────────────────────────────────────────
\section{Integración Event-Driven con la API de Groq}
% ─────────────────────────────────────────────────────────────────────────────

\subsection{Eliminación del Polling Periódico}

La versión anterior del sistema invocaba la IA de Groq cada vez que se detectaban 3 strikes de anomalía, sin distinción del tipo de cambio ni del contexto histórico. En la versión 2.0, la lógica de invocación requiere el cumplimiento simultáneo de dos condiciones:

\begin{enumerate}[noitemsep]
    \item El estado clasificado localmente es \texttt{Anomalía} o \texttt{Fuga} (no \texttt{Normal}).
    \item La función \texttt{BaselineService.esDesviacionSignificativa()} retorna \texttt{true}.
\end{enumerate}

Adicionalmente, en el caso de \texttt{Fuga}, la invocación es directa (sin esperar los 3 strikes) dado el riesgo inmediato que representa un caudal crítico.

\subsection{Estructura del Payload Contextual}

El payload enviado a Groq en la versión 2.0 incluye los siguientes campos:

\begin{center}
\begin{tabularx}{0.9\textwidth}{lX}
    \toprule
    \textbf{Campo del contexto} & \textbf{Contenido} \\
    \midrule
    Lectura anómala actual & Caudal (L/min), ruido (ADC), clasificación edge, marca temporal. \\
    Resúmenes de 5 días    & Hasta 3 resúmenes recientes de \texttt{lecturas\_resumen\_5d} en formato texto compacto. \\
    Historial mensual      & Último registro de \texttt{historial\_mensual} disponible. \\
    Baseline actual        & Promedio de flujo y ruido del período activo (día/noche). \\
    Parámetros de referencia & Umbrales del sistema (incluidos en el prompt del sistema). \\
    \bottomrule
\end{tabularx}
\captionof{table}{Campos del payload enviado a la API de Groq.}
\end{center}

\subsection{Respuesta y Normalización}

La API de Groq devuelve un JSON con tres campos (\texttt{veredicto}, \texttt{severidad}, \texttt{mensaje}). El servicio \texttt{GroqApiService} normaliza los valores recibidos a los canónicos del sistema mediante las funciones \texttt{\_normalizarVeredicto()} y \texttt{\_normalizarSeveridad()}, garantizando la consistencia semántica con la enumeración \texttt{EstadoSensor}.

En caso de fallo de red o error de la API, el sistema utiliza la clasificación local del edge como fallback, registrando el evento con el mensaje ``Clasificación edge (sin conexión IA)''. Esto garantiza que ningún evento de anomalía o fuga quede sin registrar en el historial SQLite.

% ─────────────────────────────────────────────────────────────────────────────
\section{Sistema de Alertas por Correo Electrónico}
% ─────────────────────────────────────────────────────────────────────────────

A partir de la versión 2.0, el sistema envía notificaciones por correo electrónico ante cualquier cambio de estado no-normal, con dos plantillas diferenciadas:

\begin{center}
\begin{tabularx}{0.9\textwidth}{lXX}
    \toprule
    \textbf{Estado detectado} & \textbf{Color del encabezado} & \textbf{Asunto del correo} \\
    \midrule
    Anomalía & Naranja (\texttt{\#FF9800}) & \textit{ALERTA PREVENTIVA: Anomalía Detectada} \\
    Fuga     & Rojo (\texttt{\#F44336})    & \textit{ALERTA CRÍTICA: Fuga Detectada} \\
    \bottomrule
\end{tabularx}
\captionof{table}{Diferenciación visual de plantillas de correo por nivel de estado.}
\end{center}

El correo preliminar se envía de forma inmediata al detectarse el cambio de estado (clasificación edge), sin esperar la confirmación de la IA. Posteriormente, si la IA confirma o modifica el estado, se puede enviar una segunda notificación con el veredicto definitivo. El servicio utiliza SMTP sobre Gmail (configurado mediante variables de entorno \texttt{SMTP\_EMAIL} y \texttt{SMTP\_PASSWORD}).

% ─────────────────────────────────────────────────────────────────────────────
\section{Patrones de Diseño y Convenciones de Arquitectura}
% ─────────────────────────────────────────────────────────────────────────────

\begin{center}
\begin{tabularx}{0.95\textwidth}{lXl}
    \toprule
    \textbf{Patrón} & \textbf{Aplicación en el sistema} & \textbf{Componente principal} \\
    \midrule
    MVC (Model-View-Controller) &
        Los modelos de datos son PODOs puros; los controladores encapsulan la lógica; las vistas solo leen el estado del controlador. &
        \texttt{DashboardController} \\
    \addlinespace
    Observer / ChangeNotifier &
        La UI se suscribe al controlador y se reconstruye automáticamente ante cada cambio de estado sin polling de interfaz. &
        \texttt{ChangeNotifier} + \texttt{ListenableBuilder} \\
    \addlinespace
    Command / Event-Driven &
        La IA solo se invoca cuando un evento local (desviación + strikes) lo justifica; no existe un timer de consulta periódica. &
        \texttt{\_invocarGroqPorEvento()} \\
    \addlinespace
    Sliding Window Filter &
        El buffer circular de muestras de ruido permite distinguir impactos aislados de patrones sostenidos mediante conteo en ventana temporal. &
        \texttt{RuidoFilterService} \\
    \addlinespace
    Strategy (Fallback) &
        Si la IA no responde, el sistema usa la clasificación local como estrategia de respaldo, garantizando continuidad operacional. &
        \texttt{GroqApiService.analizarConContexto()} \\
    \bottomrule
\end{tabularx}
\captionof{table}{Patrones de diseño aplicados en EdgeLeak AI v2.0.}
\end{center}

% ─────────────────────────────────────────────────────────────────────────────
\section{Glosario de Términos Técnicos}
% ─────────────────────────────────────────────────────────────────────────────

\begin{description}[style=nextline, leftmargin=3cm, font=\bfseries]
    \item[ADC] \textit{Analog-to-Digital Converter}. Valor numérico que representa la amplitud analógica del micrófono KY-037 convertida a escala digital (rango 0--4095 en ESP32 con ADC de 12 bits).
    \item[Baseline] Perfil estadístico del comportamiento normal del usuario, calculado sobre lecturas históricas de 5 días. Sirve como referencia para la detección de desviaciones.
    \item[Edge Computing] Paradigma de procesamiento en el que los datos se analizan localmente (en el dispositivo o la aplicación) antes de ser enviados a la nube, reduciendo latencia y tráfico de red.
    \item[ESP32] Microcontrolador con Wi-Fi integrado utilizado como nodo sensor. Envía las lecturas del YF-S201 y el KY-037 al servidor HTTP local de la aplicación Flutter.
    \item[Event-Driven] Modelo de invocación en el que un servicio se activa únicamente ante la ocurrencia de un evento específico (desviación del baseline), en contraposición al polling periódico.
    \item[KY-037] Módulo micrófono de alta sensibilidad con salida analógica. Utilizado para detectar patrones acústicos de goteo de agua.
    \item[L/min] Litros por minuto. Unidad de medida del caudal volumétrico de agua procesada a partir de los pulsos del sensor YF-S201.
    \item[Polling] Estrategia de consulta periódica a un servicio externo en intervalos fijos, independientemente de si ha ocurrido algún evento relevante. Reemplazada en v2.0 por el modelo event-driven.
    \item[Strike] Conteo interno de lecturas consecutivas en estado no-normal. Al alcanzar 3 strikes, se dispara la invocación de la IA (excepto en el estado \texttt{Fuga}, que es inmediato).
    \item[YF-S201] Sensor de flujo de agua tipo paleta. Genera pulsos eléctricos proporcionales al caudal volumétrico que atraviesa el sensor.
\end{description}

% ─────────────────────────────────────────────────────────────────────────────
\section{Control de Versiones del Documento}
% ─────────────────────────────────────────────────────────────────────────────

\begin{center}
\begin{tabularx}{0.9\textwidth}{lllX}
    \toprule
    \textbf{Versión} & \textbf{Fecha} & \textbf{Responsable} & \textbf{Descripción del cambio} \\
    \midrule
    1.0 & 2024-10 & PICUR Uniremington & Documento inicial. Arquitectura básica ESP32-Flutter-Groq. \\
    2.0 & 2026-05 & PICUR Uniremington &
        Incorporación de: máquina de estados de 3 niveles, filtro digital de ruido, baseline adaptativo, agregación jerárquica SQLite, invocación event-driven de IA, y correos diferenciados por severidad. \\
    \bottomrule
\end{tabularx}
\captionof{table}{Historial de versiones de la especificación técnica.}
\end{center}

\end{document}
